\chapter{Inversi\'on constitutiva y curva modular cargada}
\label{rm:chap:curva-modular-reconstruccion}

\section{La completación espectral de \(3\) a \(4\)}

El operador de vac\'io del \cref{rm:ch:vacuum} se obtiene aplicando al operador
angular la funci\'on
\begin{equation}
 F(q):=\frac{1-q^{90}}{1-q^{120}},
 \qquad 0<q<1.
 \label{rm:eq:curve-Fq}
\end{equation}
La factorizaci\'on de los dos canales por el modo elemental \(30\),
\begin{equation}
 90=3\cdot30,
 \qquad
 120=4\cdot30,
 \label{rm:eq:curve-90-120}
\end{equation}
permite introducir \(s=q^{30}\) y escribir
\begin{equation}
 F(q)=f(s),
 \qquad
 f(s):=\frac{1+s+s^2}{1+s+s^2+s^3}.
 \label{rm:eq:curve-fs}
\end{equation}
El numerador contiene tres estadios consecutivos del modo, y el denominador
incorpora el cuarto. El defecto de completaci\'on es
\begin{equation}
 1-f(s)=\frac{s^3}{1+s+s^2+s^3}>0.
 \label{rm:eq:curve-completion-defect}
\end{equation}
As\'i, la cuarta etapa queda registrada como memoria positiva dentro de la
respuesta.

\begin{theorem}[Biyectividad de la respuesta \(3\to4\)]
\label{rm:thm:curve-F-bijection}
La funci\'on \(F\) es un difeomorfismo anal\'itico decreciente
\begin{equation}
 \boxed{F:(0,1)\xrightarrow{\;\sim\;}(3/4,1).}
 \label{rm:eq:curve-F-range}
\end{equation}
Para cada \(r\in(3/4,1)\), el valor \(q=F^{-1}(r)\) se obtiene como
\(q=s^{1/30}\), donde \(s\in(0,1)\) es la ra\'iz \emph{\'unica} de
\begin{equation}
 r s^3+(r-1)(1+s+s^2)=0.
 \label{rm:eq:curve-inverse-cubic}
\end{equation}
\end{theorem}

\begin{proof}
La derivada respecto de \(s\) es
\begin{equation}
 f'(s)
 =-\frac{s^2(3+2s+s^2)}{(1+s+s^2+s^3)^2}<0
 \qquad(0<s<1).
 \label{rm:eq:curve-f-derivative}
\end{equation}
Como \(\dd s/\dd q=30q^{29}>0\), tambi\'en \(F'(q)<0\). Los l\'imites son
\[
 \lim_{q\to0^+}F(q)=1,
 \qquad
 \lim_{q\to1^-}F(q)=\frac{90}{120}=\frac34.
\]
La continuidad y la monoton\'ia estricta prueban la biyectividad; la derivada
no se anula en el intervalo, de modo que el teorema de la funci\'on inversa
prueba la analiticidad. Multiplicar
\(r=(1+s+s^2)/(1+s+s^2+s^3)\) por el denominador produce
\cref{rm:eq:curve-inverse-cubic}; la unicidad de su ra\'iz en \((0,1)\) es
equivalente a la inyectividad ya demostrada.
\end{proof}

El inversor \cref{rm:eq:curve-inverse-cubic} proporciona un procedimiento
num\'erico intr\'inseco. Una bisecci\'on en \((0,1)\) conserva el intervalo y
certifica cada d\'igito; un m\'etodo de Newton puede acelerarla usando
\cref{rm:eq:curve-f-derivative}. La evaluaci\'on decimal interviene despu\'es de
fijar la ecuaci\'on y su ra\'iz \'unica.

\section{Reconstrucci\'on exacta del estado angular desde el vac\'io}

Sean
\begin{equation}
 T=q_+P_++q_-P_-,
 \qquad
 \mathcal V=F(T)=r_+P_++r_-P_-,
 \qquad
 r_\pm=F(q_\pm).
 \label{rm:eq:curve-T-V}
\end{equation}
Las dos constitutivas adimensionales est\'an asignadas a las hojas mediante
\begin{equation}
 \widehat\varepsilon=r_+^2,
 \qquad
 \widehat\mu=r_-^2.
 \label{rm:eq:curve-mu-epsilon}
\end{equation}

\begin{theorem}[Reconstrucci\'on constitutiva reversible]
\label{rm:thm:curve-vacuum-tomography}
La respuesta constitutiva completa determina el antecedente angular. En la
base espectral etiquetada,
\begin{align}
 \boxed{q_+=F^{-1}(\sqrt{\widehat\varepsilon}),}
 &\qquad
 \boxed{q_-=F^{-1}(\sqrt{\widehat\mu}),}
 \label{rm:eq:curve-recover-qpm}\\
 \boxed{D_A=q_+q_-}
 &=F^{-1}(\sqrt{\widehat\varepsilon})
   F^{-1}(\sqrt{\widehat\mu}),
 \label{rm:eq:curve-recover-DA}\\
 \boxed{A_{\rm rad}=-\frac12\log D_A,}
 &\qquad
 \boxed{C^*_{\rm rad}=\frac12\log\frac{q_-}{q_+}.}
 \label{rm:eq:curve-recover-A-C}
\end{align}
Si se conoce el operador \(\mathcal V\), en lugar de sus dos autovalores
etiquetados, la reconstrucci\'on adopta la forma de c\'alculo funcional
\begin{equation}
 \boxed{T=F^{-1}(\mathcal V).}
 \label{rm:eq:curve-functional-inverse}
\end{equation}
\end{theorem}

\begin{proof}
Los autovalores de \(\mathcal V\) pertenecen a \((3/4,1)\) y la ra\'iz
positiva de cada constitutiva es \(r_\pm\). El
\cref{rm:thm:curve-F-bijection} proporciona un \'unico \(q_\pm\) para cada
hoja. Las f\'ormulas de \(D_A,A_{\rm rad},C^*_{\rm rad}\) son las inversas
de \cref{rm:eq:modular-determinant-ratio}. Para la forma operatoria, el teorema
espectral aplica la funci\'on continua \(F^{-1}\) al espectro de
\(\mathcal V\).
\end{proof}

Cuando s\'olo se conservan los dos n\'umeros
\((\widehat\mu,\widehat\varepsilon)\), la asignaci\'on a hojas forma parte
de los datos tipados. El operador \(\mathcal V\) contiene adem\'as sus
proyectores espectrales y reconstruye la representaci\'on completa. Esta
distinci\'on separa la reconstrucci\'on del espectro de la elecci\'on de una
base.

Las restantes lecturas se recuperan simult\'aneamente:
\begin{equation}
 \widehat Z=\frac{r_-}{r_+}
 =\sqrt{\frac{\widehat\mu}{\widehat\varepsilon}},
 \qquad
 \widehat c=\frac1{r_-r_+}
 =\frac1{\sqrt{\widehat\mu\widehat\varepsilon}}.
 \label{rm:eq:curve-recover-Z-c}
\end{equation}
El producto de hojas publica la propagaci\'on com\'un; el cociente publica la
orientaci\'on constitutiva. Bajo \(C^*\mapsto-C^*\), las dos constitutivas se
intercambian, \(\widehat Z\) se invierte y \(\widehat c\) permanece fija.

\begin{corollary}[Control cruzado electrod\'ebil]
\label{rm:cor:curve-weak-control}
En la interfaz angular donde
\((m_W^{\angle}/m_Z^{\angle})^2=D_A\), la reconstrucci\'on del vac\'io
implica
\begin{equation}
 \boxed{
 \sin^2\theta_W^{\rm OS,HMT}
 =1-F^{-1}(\sqrt{\widehat\varepsilon})
      F^{-1}(\sqrt{\widehat\mu}).}
 \label{rm:eq:curve-weak-from-vacuum}
\end{equation}
\end{corollary}

\begin{proof}
El \cref{rm:thm:curve-vacuum-tomography} reconstruye \(D_A\). La interfaz
angular declarada define
\(\sin^2\theta_W^{\rm OS}=1-(m_W/m_Z)^2=1-D_A\).
\end{proof}

La ecuaci\'on \cref{rm:eq:curve-weak-from-vacuum} es un control entre lectoras
hermanas. La genealog\'ia causal contin\'ua partiendo de \(T\): el vac\'io
aplica \(F(T)\), y el lector electrod\'ebil aplica \(\det T\). La
invertibilidad de \(F\) permite comparar ambas publicaciones sin convertir
una de ellas en fundamento de la otra.

\section{La curva modular cargada}

La identidad
\begin{equation}
 D_A=\exp\!\left(-\frac{100\pi}{9}\alpha_{\HMT}\right)
 \label{rm:eq:curve-D-alpha}
\end{equation}
sugiere usar el determinante como coordenada global de la familia. Para
\(D\in(0,1)\), definimos
\begin{equation}
 \boxed{
 Q(D):=-\frac9{25}\log D,}
 \qquad
 a(D):=\frac{Q(D)}{4\pi}
 =-\frac9{100\pi}\log D.
 \label{rm:eq:curve-Q-a}
\end{equation}
En el punto HMT, \(a(D_A)=\alpha_{\HMT}\) y
\begin{equation}
 \boxed{Q(D_A)=4\pi\alpha_{\HMT}=:e_g^2.}
 \label{rm:eq:curve-Q-gauge}
\end{equation}
El s\'imbolo \(Q(D)\) designa aqu\'i el acoplamiento cargado adimensional;
se distingue de cualquier secci\'on dimensional de carga o energ\'ia.

Para que el lector regional sea autocontenido, introducimos el car\'acter
qu\'intico
\begin{equation}
 \begin{aligned}
 H_5(a):={}&\frac12\sqrt{\frac{1000a}{\varphi_{\HMT}}}-a
 +\frac9{16}a^2-\frac59a^3
 +\frac7{48}a^4-\frac1{54}a^5,
 \end{aligned}
 \label{rm:eq:curve-H5}
\end{equation}
el germen regional
\begin{equation}
 \mathcal C_\pi(a,D)
 :=169a^6+\frac{Da^7}{1-Da},
 \label{rm:eq:curve-Cpi}
\end{equation}
y la fase orientada
\begin{equation}
 y(D):=\frac{\pi}{90}\bigl(H_5(a(D))+a(D)\bigr)
 -\mathcal C_\pi(a(D),D).
 \label{rm:eq:curve-yD}
\end{equation}
El selector \(169\), el t\'ermino finito de grado seis y la prolongaci\'on
geom\'etrica de grado siete conservan su procedencia regional. La condici\'on
\(0<Da(D)<1\) garantiza la convergencia del germen.

Definimos el cociente orientado
\begin{equation}
 \varrho_{\rm or}(D):=\ee^{2y(D)}
 \label{rm:eq:curve-rho-or}
\end{equation}
y el intervalo admisible
\begin{equation}
 \mathcal I_{\rm mod}:=
 \left\{D\in(0,1):
 0<Da(D)<1,
 \ |y(D)|<-\frac12\log D
 \right\}.
 \label{rm:eq:curve-domain}
\end{equation}
La segunda desigualdad asegura que los dos canales reconstruidos son
contracciones estrictas.

\begin{definition}[Curva modular cargada]
\label{rm:def:charged-modular-curve}
La curva modular cargada es la aplicaci\'on
\begin{equation}
 \Gamma_{\rm ch}:\mathcal I_{\rm mod}\longrightarrow(0,1)\times\R_{>0},
 \qquad
 \boxed{\Gamma_{\rm ch}(D)=\bigl(D,\varrho_{\rm or}(D)\bigr).}
 \label{rm:eq:curve-Gamma}
\end{equation}
Sobre ella,
\begin{equation}
 q_-(D)=\sqrt{D\varrho_{\rm or}(D)},
 \qquad
 q_+(D)=\sqrt{\frac{D}{\varrho_{\rm or}(D)}}.
 \label{rm:eq:curve-qD}
\end{equation}
\end{definition}

\begin{theorem}[Rango funcional uno]
\label{rm:thm:curve-rank-one}
La curva \(\Gamma_{\rm ch}\) es una inmersi\'on anal\'itica de rango uno.
El espacio gen\'erico de contracciones positivas bicapa requiere dos
coordenadas independientes \((D,\varrho_{\rm or})\); el lector HMT restringe
ese espacio a la gr\'afica \(\varrho_{\rm or}=\varrho_{\rm or}(D)\).
Adem\'as,
\begin{equation}
 A_{\rm rad}(D)=-\frac12\log D
 =\frac{50\pi}{9}a(D),
 \qquad
 C^*_{\rm rad}(D)=y(D).
 \label{rm:eq:curve-A-C-D}
\end{equation}
\end{theorem}

\begin{proof}
Las funciones de \cref{rm:eq:curve-Q-a,rm:eq:curve-H5,rm:eq:curve-Cpi,rm:eq:curve-yD}
son anal\'iticas en \(\mathcal I_{\rm mod}\). La primera componente de
\(\Gamma_{\rm ch}'(D)\) es uno, luego la derivada nunca se anula y su rango
es uno. Las identidades \cref{rm:eq:curve-A-C-D} siguen directamente de las
definiciones. Finalmente, el producto y el cociente de
\cref{rm:eq:curve-qD} son \(D\) y \(\varrho_{\rm or}(D)\), respectivamente.
\end{proof}

La familia resolvente del cap\'itulo anterior se restringe a la misma curva:
\begin{equation}
 \mathfrak M_D(z)
 :=Q(D)^{-1}\bigl(zI-T(D)\bigr)^{-1},
 \qquad
 T(D)=q_+(D)P_++q_-(D)P_-.
 \label{rm:eq:curve-resolvent-D}
\end{equation}
En el punto \(D_A\), la igualdad \(Q(D_A)=4\pi\mathfrak A_\beta\)
identifica \(\mathfrak M_D\) con \(\mathfrak M_\beta\). El acoplamiento,
la intensidad, la orientaci\'on y el vac\'io quedan as\'i coordinados por una
misma variable geneal\'ogica, mientras sus funciones lectoras permanecen
tipadas.

\section{Confluencia entre vac\'io, Higgs e impedancia}

El lector del vac\'io publica a lo largo de la curva
\begin{align}
 \widehat\mu(D)&=F(q_-(D))^2,
 &\widehat\varepsilon(D)&=F(q_+(D))^2,
 \label{rm:eq:curve-vacuum-D}\\
 \widehat Z(D)&=\frac{F(q_-(D))}{F(q_+(D))},
 &\widehat c(D)&=\frac1{F(q_-(D))F(q_+(D))}.
 \label{rm:eq:curve-Z-c-D}
\end{align}
Por la inyectividad de \(F\), cualquiera de las parejas operatorias
completas reconstruye \((D,\varrho_{\rm or})\).

La ruta angular de Higgs contiene la orientaci\'on mediante
\begin{equation}
 \frac{m_H^{\angle}}{m_W^{\angle}}
 =\exp\!\left(3A_{\rm rad}+\frac53C^*_{\rm rad}\right).
 \label{rm:eq:curve-HW-ratio}
\end{equation}
En la interfaz de \'arbol con
\(m_H^2=2\lambda_Hv^2\) y \(m_W^2=g^2v^2/4\), esta identidad equivale a
\begin{equation}
 \lambda_H
 =\frac{Q(D)}{8(1-D)}D^{-3}
  \varrho_{\rm or}(D)^{5/3}.
 \label{rm:eq:curve-lambda-H}
\end{equation}
Por tanto, el escalar
\begin{equation}
 \Xi_H(D,\lambda_H)
 :=\frac{8(1-D)D^3\lambda_H}{Q(D)}
 \label{rm:eq:curve-Xi-H}
\end{equation}
satisface
\begin{equation}
 \boxed{
 \Xi_H=\varrho_{\rm or}^{5/3},
 \qquad
 \varrho_{\rm or}=\Xi_H^{3/5}.}
 \label{rm:eq:curve-Higgs-rho}
\end{equation}

\begin{corollary}[Reconstrucci\'on Higgs--impedancia]
\label{rm:cor:curve-Higgs-impedance}
Bajo la interfaz de \'arbol declarada y para \(\Xi_H>0\),
\begin{equation}
 \boxed{
 \widehat Z
 =\frac{
 F\!\left(\sqrt D\,\Xi_H^{3/10}\right)}
 {F\!\left(\sqrt D\,\Xi_H^{-3/10}\right)}.}
 \label{rm:eq:curve-Higgs-Z}
\end{equation}
El mismo cociente orientado se reconstruye por tres v\'ias:
\begin{equation}
 \varrho_{\rm or}^{\rm HMT}(D)
 =\frac{F^{-1}(\sqrt{\widehat\mu})}
        {F^{-1}(\sqrt{\widehat\varepsilon})}
 =\Xi_H(D,\lambda_H)^{3/5}.
 \label{rm:eq:curve-triple-confluence}
\end{equation}
\end{corollary}

\begin{proof}
La ecuaci\'on \cref{rm:eq:curve-Higgs-rho} da
\(q_-=\sqrt D\,\Xi_H^{3/10}\) y
\(q_+=\sqrt D\,\Xi_H^{-3/10}\). Sustituirlos en
\(\widehat Z=F(q_-)/F(q_+)\) prueba \cref{rm:eq:curve-Higgs-Z}. La segunda
igualdad de \cref{rm:eq:curve-triple-confluence} es
\cref{rm:eq:curve-Higgs-rho}; la primera procede del
\cref{rm:thm:curve-vacuum-tomography}.
\end{proof}

La sustituci\'on de \(\lambda_H\) obtenida de
\cref{rm:eq:curve-lambda-H} convierte
\cref{rm:eq:curve-triple-confluence} en una identidad algebraica de consistencia.
Una determinaci\'on de \(\lambda_H\) desde un lector independiente transforma
la misma ecuaci\'on en una prueba cruzada falsable de la orientaci\'on. Esta
separaci\'on evita atribuir independencia probatoria a una mera inversi\'on de
la ecuaci\'on de partida.

La repartici\'on de calibre adopta tambi\'en una forma puramente determinada
por \(D\):
\begin{equation}
 g^2=\frac{Q(D)}{1-D},
 \qquad
 g'^2=\frac{Q(D)}D,
 \qquad
 \boxed{
 \frac1{g^2}+\frac1{g'^2}=\frac1{Q(D)}
 =-\frac{25}{9\log D}.}
 \label{rm:eq:curve-gauge-harmonic}
\end{equation}
El acoplamiento total \(Q(D)\) se reparte entre los dos sectores de calibre,
y la suma arm\'onica recupera exactamente el macroobjeto cargado.

\section{Empalme de secciones causales compatible con orientación y datos de frontera}

Conviene distinguir el acoplamiento \(Q(D)\) de la energ\'ia de Planck,
que denotaremos por \(\mathcal Q_P=E_P\). La pareja de Planck define una
velocidad interna
\begin{equation}
 c_{\rm int}:=\frac{\mathcal Q_P\ell_P}{\hbar},
 \label{rm:eq:curve-c-int}
\end{equation}
mientras el operador constitutivo define
\begin{equation}
 c_{\rm vac}:=\frac1{\sqrt{\mu_{\rm vac}\varepsilon_{\rm vac}}}.
 \label{rm:eq:curve-c-vac}
\end{equation}
Ambas expresiones viven en rectas causales tipadas. Un morfismo de costura
\begin{equation}
 J_C:L_C^{\rm vac}\longrightarrow L_C^{\rm grav}
 \label{rm:eq:curve-JC}
\end{equation}
permite compararlas, y su defecto adimensional es
\begin{equation}
 \boxed{
 \Xi_c
 :=\frac{\mathcal Q_P\ell_P}{\hbar}
   \sqrt{\mu_{\rm vac}\varepsilon_{\rm vac}}
 =\frac{c_{\rm int}}{c_{\rm vac}}.}
 \label{rm:eq:curve-Xi-c}
\end{equation}

\begin{criterion}[Costura causal]
\label{rm:crit:curve-causal-seam}
El vac\'io electromagn\'etico y la realizaci\'on gravitatoria publican una
misma secci\'on causal exactamente cuando
\begin{equation}
 \boxed{\Xi_c=1,}
 \qquad
 J_C(c_{\rm vac})=c_{\rm int}.
 \label{rm:eq:curve-Xi-one}
\end{equation}
\end{criterion}

Una vez satisfecha esta costura, la identidad CT108
\(
G=(54/\pi)^2c^5t_0^2/\hbar
\)
se expresa enteramente mediante el producto constitutivo:
\begin{equation}
 \boxed{
 G=\left(\frac{54}{\pi}\right)^2
 \frac{t_0^2}
 {\hbar(\mu_{\rm vac}\varepsilon_{\rm vac})^{5/2}}.}
 \label{rm:eq:curve-G-from-vacuum}
\end{equation}
La gravedad emplea el producto par de las hojas; la impedancia y la carga
emplean el cociente orientado
\begin{equation}
 Z_{\rm vac}^2=\frac{\mu_{\rm vac}}{\varepsilon_{\rm vac}}.
 \label{rm:eq:curve-Z-oriented}
\end{equation}
La pareja producto--cociente distribuye, por tanto, causalidad y orientaci\'on
sin disgregar el antecedente espectral.

\section{Reconstrucci\'on estable en la torre UHF}

Sea \(\tau_m=\tau_2\otimes\tau_{9^m}\) la traza normalizada de
\(\mathcal A_m=M_2\otimes M_{9^m}\), y sean
\begin{equation}
 T_m=T\otimes I_{9^m},
 \qquad
 \mathcal V_m=F(T_m)=\mathcal V\otimes I_{9^m}.
 \label{rm:eq:curve-UHF-TV}
\end{equation}
La inversi\'on funcional conmuta con el refinamiento:
\begin{equation}
 F^{-1}(\mathcal V_m)=T_m,
 \qquad
 E_m(\mathcal V_{m+1})=\mathcal V_m,
 \qquad
 E_m(T_{m+1})=T_m.
 \label{rm:eq:curve-UHF-inversion}
\end{equation}

El determinante matricial ordinario de \(T_m\) es \(D_A^{9^m}\). El
invariante compatible con el refinamiento es el determinante logar\'itmico
normalizado
\begin{equation}
 \boxed{
 \mathfrak d_m(T_m)
 :=\exp\!\left(2\tau_m(\log T_m)\right)=D_A.}
 \label{rm:eq:curve-normalized-determinant}
\end{equation}
La orientaci\'on posee la prolongaci\'on igualmente estable
\begin{equation}
 \boxed{
 \mathfrak r_m(T_m,R_m)
 :=\exp\!\left(-2\tau_m(R_m\log T_m)\right)
 =\varrho_{\rm or},}
 \qquad
 R_m=R\otimes I_{9^m}.
 \label{rm:eq:curve-normalized-orientation}
\end{equation}

\begin{proposition}[Estabilidad de la reconstrucci\'on a toda profundidad]
\label{rm:prop:curve-UHF-tomography}
Para todo \(m\), el par
\begin{equation}
 \left(
 \mathfrak d_m(F^{-1}(\mathcal V_m)),
 \mathfrak r_m(F^{-1}(\mathcal V_m),R_m)
 \right)
 \label{rm:eq:curve-UHF-pair}
\end{equation}
es igual a \((D_A,\varrho_{\rm or})\). En consecuencia, reconstruye las
mismas coordenadas \(A_{\rm rad}\) y \(C^*_{\rm rad}\) en todos los
niveles y en el l\'imite UHF.
\end{proposition}

\begin{proof}
Como
\(\log T_m=(\log T)\otimes I_{9^m}\), la traza normalizada elimina el
factor de amplificaci\'on. Adem\'as,
\[
 \tau_2(\log T)=\frac12(\log q_++\log q_-)=\frac12\log D_A
\]
y
\[
 \tau_2(R\log T)=-y=-\frac12\log\varrho_{\rm or}.
\]
Las ecuaciones \cref{rm:eq:curve-normalized-determinant,rm:eq:curve-normalized-orientation}
siguen al exponentiar. La inversi\'on funcional de
\cref{rm:eq:curve-UHF-inversion} completa la reconstrucci\'on.
\end{proof}

Esta formulaci\'on muestra que la completaci\'on \(3\to4\) conserva su
capacidad reconstructiva a toda profundidad. Los cilindros non\'adicos
aumentan la resoluci\'on del estado; el determinante normalizado, la
orientaci\'on, la respuesta y la curva modular permanecen naturales respecto
de las expectativas.

\section{Separación entre identidades cuánticas internas, interfaces físicas y obstrucciones de la reconstrucción}

\begin{proofstatusbox}
La monoton\'ia y la biyectividad de \(F\), la inversi\'on operatoria, la
reconstrucci\'on \((q_+,q_-,D_A,C^*)\), el rango uno de
\(\Gamma_{\rm ch}\), la identidad arm\'onica de calibre y la estabilidad UHF
son resultados exactos. La curva modular cargada formaliza una arquitectura
autoral preexistente con estructura de partida expl\'icita: cierre de
\(\alpha_{\HMT}\), selector regional \(169\), car\'acter \(H_5\), germen
\(\mathcal C_\pi\) y orientaci\'on. La igualdad
Higgs--impedancia es exacta dentro de la interfaz de \'arbol; constituye una
prueba independiente solamente cuando \(\lambda_H\) procede de un lector
distinto. La condici\'on \(\Xi_c=1\) y la ecuaci\'on gravitatoria posterior
son una interfaz f\'isica tipada cuya promoci\'on exige el morfismo \(J_C\).
\end{proofstatusbox}

\begin{falsifierbox}
La reconstrucci\'on queda refutada si \(F'(q)\) cambia de signo en
\((0,1)\), si la ecuaci\'on c\'ubica admite dos ra\'ices en ese intervalo, si
\(F^{-1}(\mathcal V)\) no recupera \(T\), o si las constitutivas reconstruyen
un producto distinto de \(D_A\). La curva queda refutada si
\(Q(D_A)\ne4\pi\alpha_{\HMT}\), si su derivada pierde rango, si
\(q_\pm(D)\) abandona \((0,1)\) dentro del dominio declarado, o si las tres
reconstrucciones de \(\varrho_{\rm or}\) en
\cref{rm:eq:curve-triple-confluence} discrepan bajo una misma normalizaci\'on.
La costura causal queda refutada por \(\Xi_c\ne1\) despu\'es de aplicar el
morfismo de rectas. En la torre, usar el determinante ordinario como si fuera
invariante producir\'ia \(D_A^{9^m}\); el control correcto exige
\cref{rm:eq:curve-normalized-determinant}.
\end{falsifierbox}
